\documentclass[12pt,a4paper]{article}
\usepackage[latin1,utf8]{inputenc}
\usepackage[portuguese]{babel}
%\usepackage[T1]{fontenc}
\usepackage{multicol}
\usepackage{amsmath}
%\usepackage{graphicx}
\usepackage{amssymb}
%\usepackage{amsfonts}
\usepackage{color}
\usepackage{multirow}
\usepackage{multicol}


\author{Autor \\ Lafaiete Santos \\ Versão \textbf{XX}}
\title{\textbf{Teste de Produção de Relatórios}}
\date{}

\begin{document}
\maketitle

\section*{\Huge Modo Matemático}
\vspace{1cm}

\begin{multicols}{2}

\subsection{Combinação}
$$ {n \choose k} = {n \choose n-k}$$

\subsection{Matriz}
$$A=\begin{pmatrix}
1+2i&&3-i&&-i\\
2-3i&&4+i&&5\\
\end{pmatrix}\bigstar$$
\end{multicols}
\vspace{1cm}

\subsection{Equações}
\begin{equation}
\begin{split}
x_1=y_1-z_1&+w_1\\
x_2=y_2&+z_2-y_3\\
\end{split}
\end{equation}
\begin{equation}
\begin{split}
x_3=&-w_3\\
x_4=y_4+z_4&-w_4-v_4\\
\end{split}
\end{equation}
\vspace{1cm}

\begin{multicols}{2}
\subsection{Expressão}
$$
\mathbb{Q}^+ = \left\{\frac{a}{b}:a,b \in\mathbb{Z}^+ \text{ e }b\neq0\right\}$$

\subsection{Sistema de Equações}
$$
y = \begin{cases} 
2x+5 & \text{se } x<1\\
a^2+b^2 & \text{se } x>0 \text{ e } y\leq2\\
0 & \text{qualquer outro valor}\\
\end{cases} $$
\end{multicols}
\vspace{1cm}

\subsection{Integral e Somatório}
$$
\int_0^{\pi} x^ne^{-ax}dx= \sum_{k=0}^{n}\dfrac{(-1)^k\cdot n!}{k!(n-k!)} \cdot \dfrac{\pi^{n-k+1}}{a^{n-k+1}} $$
\newpage

\section*{\Huge Modo Texto}
\subsubsection{Texto e Formatação com espaço de 16 milímetros}
\hskip 16mm A Universidade de {\large Lisboa} recebe todos os anos cerca de \textcolor{blue}{9.000 estudantes} internacionais provenientes de mais de \textcolor{red}{100 países} que procuram uma educação de alta qualidade, bem como a cultura, o clima e a hospitalidade que Lisboa, e {\Huge Portugal}, têm para oferecer. A ULisboa é a
\underline{maior Universidade de Portugal} e uma das maiores da {\Huge Europa}.\footnote{Texto retirado de https://ciencias.ulisboa.pt/}

\subsubsection{Tabela Alinhada À Direita}
\begin{table}[h!]
\begin{flushright}
\begin{tabular}{|c||l||r|}\hline
Total Estudantes & Grau de Estudo & Nº de Estudantes \\ \hline \hline
\multirow{4}{*}
{\textbf{51.917}} & Licenciatura & \textbf{26.326} \\
& Mestrado Integrado & \textbf{6.419} \\ 
& Mestrado & \textbf{14.212} \\ 
& Doutoramento & \textbf{4.961} \\ \hline
\end{tabular}
\caption{Estudantes na ULisboa}
\end{flushright}
\end{table}

\begin{multicols}{2}

\subsubsection{Lista}

\subsubsection{Outros}


\end{multicols}

\end{document}